\documentclass[11pt]{amsart}
\usepackage{amsmath,amssymb,amsthm,booktabs}
\usepackage[margin=1.15in]{geometry}
\usepackage[hidelinks]{hyperref}
\newtheorem{theorem}{Theorem}[section]
\newtheorem{lemma}[theorem]{Lemma}
\theoremstyle{remark}
\newtheorem{remark}[theorem]{Remark}
\newcommand{\PP}{\mathbb{P}}
\newcommand{\GP}{\mathrm{GP}_b}
\newcommand{\doi}[1]{\href{https://doi.org/#1}{doi:\nolinkurl{#1}}}
\newcommand{\arxiv}[1]{\href{https://arxiv.org/abs/#1}{arXiv:\nolinkurl{#1}}}
\setlength{\emergencystretch}{2.5em}

\title[Cluster repulsion on grandparent graphs]{Finite touching sets between percolation clusters on grandparent
graphs}
\author{Yitzchak Shmalo}
\address{Einstein Institute of Mathematics, The Hebrew University of Jerusalem, Jerusalem, Israel}
\email{yspennstate@gmail.com}
\date{November 2026}
\subjclass[2020]{60K35, 82B43, 05C63}
\keywords{Bernoulli percolation, cluster repulsion, nonunimodular graphs, grandparent graph, two-point function}

\raggedbottom
\begin{document}
\begin{abstract}
A transitive graph exhibits cluster repulsion at the parameter $p$ if, almost surely, any two distinct clusters of
Bernoulli bond percolation with parameter $p$ are joined by only finitely many edges. Timár proved cluster repulsion for
every unimodular quasi-transitive graph and every $p$, and the nonunimodular case is open. Bell, Chu, Rodgers, Terlov and
Tserunyan proved that on nonunimodular graphs the touching sets have finite modular weight, which leaves open whether
two light clusters can touch along infinitely many edges. We prove cluster repulsion at every $p\in[0,1]$ for the
grandparent graphs $\GP$, $b\ge2$, the standard transitive nonunimodular examples, whose infinite clusters are light for
$1/b<p<1$. The proof shows that every two-point function $p\mapsto\PP_p(u\leftrightarrow v)$ is continuous on $(0,1)$,
by integrating out each descendant cone through an exact one-variable equation and bounding the effect of the missing
ancestors uniformly in $p$, and then applies the sprinkling argument of Häggström, Peres and Schonmann.
\end{abstract}
\maketitle

\section{Introduction}\label{sec:intro}

Let $G=(V,E)$ be an infinite, connected, locally finite transitive graph, and let $\PP_p$ be the law of Bernoulli bond
percolation on $G$ with parameter $p\in[0,1]$: each edge is open with probability $p$, independently. For distinct open
clusters $C$ and $D$ the touching set is
\[
\tau(C,D)=\{x\in C:\ \{x,y\}\in E\text{ for some }y\in D\}.
\]
Since $G$ has bounded degree, $\tau(C,D)$ is finite if and only if finitely many edges join $C$ and $D$. The graph
exhibits cluster repulsion at $p$ if, $\PP_p$-almost surely, $\tau(C,D)$ is finite for every pair of distinct clusters.

Häggström, Peres and Schonmann asked whether every quasi-transitive graph exhibits cluster repulsion for every
$p$~\cite[p.~87]{HPS99}, and observed that it holds at every $p$ at which all two-point functions
$p\mapsto\PP_p(u\leftrightarrow v)$ are continuous. Timár~\cite[Theorem~1.1]{Tim06a} proved cluster repulsion for every
unimodular quasi-transitive graph. The survey of Häggström and Jonasson~\cite[end of Section~5]{HJ06} records the general
statement as a conjecture and notes that the nonunimodular case remains open.

On a nonunimodular graph, mass transport must be weighted by the modular function, and an infinite set can have finite
total weight. Fix a closed transitive group $\Gamma$ of automorphisms with left Haar measure $m$ and a vertex $o$, and
give $x$ the weight $m(\Gamma_x)/m(\Gamma_o)$, where $\Gamma_x$ is the stabilizer of $x$. A set of vertices is heavy if
its total weight is infinite and light otherwise. Bell, Chu, Rodgers, Terlov and Tserunyan~\cite[Theorem~1.2]{BCRTT}
proved that on every quasi-transitive graph the touching set of two distinct infinite clusters has finite weight. Finite
weight gives finite cardinality only when the weights are bounded away from zero, which fails on a nonunimodular graph:
on the grandparent graph of Section~\ref{sec:gp}, a downward ray has weights $1,b^{-1},b^{-2},\dots$. On a nonunimodular
transitive graph a heavy infinite cluster and a light infinite cluster do not coexist, by the robust-property theorem
of~\cite[Theorem~4.1.6]{HPS99} (see also~\cite{Tang19}). So what remains open of the question of Häggström, Peres and
Schonmann~\cite[Question~1.1]{BCRTT} is whether two light infinite clusters can touch along infinitely many edges. Light infinite clusters occur on every
nonunimodular transitive graph, since $p_c<p_h$ there~\cite{Hut20}.

\subsection{Grandparent graphs}\label{sec:gp}
Let $b\ge2$, let $T$ be the $(b+1)$-regular tree, fix an end $\xi$ of $T$, and for a vertex $x$ let $P(x)$ be its
neighbour in the direction of $\xi$. The grandparent graph $\GP$ has the vertex set of $T$, the edges of $T$, and the
edges $\{x,P^2(x)\}$ for every vertex $x$. It is transitive and nonunimodular (Lemma~\ref{lem:structure}), and it is the
standard example of such a graph; percolation on it is discussed in \cite{HPS99,Tim06b,Spa09,MS26}. For $1/b<p<1$ it has
infinitely many infinite clusters, all light (Remark~\ref{rem:light}).

\begin{theorem}\label{thm:main}
Let $b\ge2$ and $p\in[0,1]$. Then $\PP_p$-almost surely every two distinct clusters of Bernoulli bond percolation on $\GP$
have finite touching sets.
\end{theorem}

The statement is for each fixed $p$, and for that $p$ it covers all pairs of clusters simultaneously. The proof does not
use the weighted repulsion theorem of~\cite{BCRTT}. The equation~\eqref{eq:cone} for the reliability of a descendant cone
is equivalent to equation~(7) of Ménard and Sénizergues~\cite{MS26} in the case of deterministic offspring, where it is
used to compute $p_c$, and the decomposition of $\GP$ into pieces attached at two vertices goes back to
Špakulová~\cite[Section~5.3]{Spa09}.

\subsection{Idea of the proof}
By the observation of Häggström, Peres and Schonmann (Lemma~\ref{lem:sprinkle}), it suffices to show that each two-point
function is continuous on $(0,1)$. Fix $u,v$ and an ancestor $a$ of both. The descendant cones that hang off the finite
tree spanned by $u$, $v$ and $a$ attach to it at no more than two vertices each, and each can be replaced by a single edge
whose probability is an explicit continuous function of $p$ (Section~\ref{sec:cone}). This gives the connection
probability inside the cone of $a$ as a continuous function of $p$. The edges that leave the cones of $a,P^3(a),P^6(a),
\dots$ form disjoint sets of fixed size, and if all edges of one of them are closed the cluster of $u$ cannot leave that
cone. This bounds the effect of the ancestors above $P^n(a)$ by a geometric term, uniformly for $p$ in compact subsets of
$(0,1)$ (Section~\ref{sec:proof}).

\section{Structure of the grandparent graph}\label{sec:structure}

Fix a vertex $o$ and let $h$ be the height function on $T$ with $h(o)=0$ and $h(P(x))=h(x)+1$. For a vertex $x$, the
children of $x$ are the $b$ vertices $y$ with $P(y)=x$, and the descendant cone $D(x)$ is the set consisting of $x$ and
all vertices $y$ with $P^k(y)=x$ for some $k\ge1$.

\begin{lemma}\label{lem:structure}
The graph $\GP$ is transitive and has degree $b^2+b+2$. Every automorphism of $\GP$ commutes with $P$, and the weights
are $m(\Gamma_x)/m(\Gamma_o)=b^{h(x)}$. In particular $\GP$ is not unimodular.
\end{lemma}

\begin{proof}
A vertex $x$ is adjacent to $P(x)$, $P^2(x)$, its $b$ children and its $b^2$ grandchildren, so its degree is $b^2+b+2$.
The automorphisms of $T$ that fix $\xi$ preserve $P$, hence preserve the added edges, and they act transitively.

The relation $P$ can be read off from the unlabelled graph. The common neighbours of the endpoints of a tree edge
$\{y,P(y)\}$ are $P^2(y)$ and the $b$ children of $y$, which is $b+1$ vertices, while the only common neighbour of the
endpoints of an added edge $\{y,P^2(y)\}$ is $P(y)$. So the tree edges are determined by the graph. Among the tree
neighbours of $x$, two are joined by an added edge exactly when one of them is $P(x)$ and the other is a child of $x$.
This relation is a star with $b\ge2$ leaves, its centre $P(x)$ is unique, and so every automorphism maps $P(x)$ to
$P(\varphi(x))$.

The stabilizer of $x$ fixes $P(x)$, while the stabilizer of $P(x)$ permutes the $b$ children of $P(x)$ transitively.
Therefore $|\Gamma_x P(x)|=1$ and $|\Gamma_{P(x)}x|=b$, and $m(\Gamma_{P(x)})/m(\Gamma_x)=|\Gamma_{P(x)}x|/|\Gamma_x
P(x)|=b$. Iterating gives the weights $b^{h(x)}$, which are not constant, so $\GP$ is not unimodular.
\end{proof}

\begin{remark}\label{rem:light}
Let $1/b<p<1$. Choose a ray $x_0,x_1=P(x_0),\dots$ and for each $k$ a child $z_k\neq x_{k-1}$ of $x_k$ (any child for
$k=0$). The cones $D(z_k)$ are disjoint. Consider the event that all edges joining $D(z_k)$ to its complement are closed
and that the open tree edges inside $D(z_k)$ contain an infinite path from $z_k$. The second condition has positive
probability because the open tree edges below $z_k$ form a Galton--Watson tree with Binomial$(b,p)$ offspring, of mean
$bp>1$. These events have the same positive probability and depend on disjoint sets of edges, so infinitely many of them
occur almost surely, and each produces an infinite cluster contained in a cone. Häggström, Peres and Schonmann record
that $p_c<p_h=p_u=1$ on $\GP$~\cite{HPS99}, so for $p<1$ there is no heavy cluster and every infinite cluster is light.
In this range Theorem~\ref{thm:main} is not a consequence of the weighted statement of~\cite{BCRTT}.
\end{remark}

\section{The reliability of a descendant cone}\label{sec:cone}

For a vertex $x$ let $H_x$ be the subgraph of $\GP$ induced by $D(x)\cup\{P(x)\}$, and let
\[
r(p)=\PP_p\bigl(x\leftrightarrow P(x)\text{ by an open path in }H_x\bigr),
\]
which does not depend on $x$.

\begin{lemma}\label{lem:cone}
For $0<p<1$, $r(p)$ is the unique solution $t\in[0,1]$ of
\begin{equation}\label{eq:cone}
t=F_p(t),\qquad F_p(t)=1-(1-p)(1-pt)^b,
\end{equation}
and $r$ is continuous on $(0,1)$.
\end{lemma}

\begin{proof}
The edges of $H_x$ at $P(x)$ are $\{x,P(x)\}$ and the edges $\{y,P(x)\}$ for the $b$ children $y$ of $x$. For a child
$y$, the edges with exactly one endpoint in $D(y)$ that lie in $H_x$ are $\{y,x\}$, $\{y,P(x)\}$ and the edges joining
the children of $y$ to $x$. Take a simple open path from $x$ to $P(x)$ in $H_x$. Its last edge is either $\{x,P(x)\}$ or
$\{y,P(x)\}$ for a child $y$. In the second case, the part of the path after its last visit to $x$ runs from $x$ to $y$
inside $D(y)\cup\{x\}$, so $y\leftrightarrow x$ in $H_y$. Hence $x\leftrightarrow P(x)$ in $H_x$ if and only if
$\{x,P(x)\}$ is open or, for some child $y$, the edge $\{y,P(x)\}$ is open and $y\leftrightarrow x$ in $H_y$. The graphs
$H_y$ for different children $y$ have disjoint edge sets, so these $b+1$ events are independent. The first has
probability $p$ and each of the others $pr(p)$, which gives $r(p)=F_p(r(p))$.

The function $g(t)=F_p(t)-t$ is concave on $[0,1]$ with $g(0)=p>0$ and $g(1)=-(1-p)^{b+1}<0$. If $0<t_1<t_2$ were
two zeros, concavity would give $g(t_1)\ge (1-t_1/t_2)g(0)+(t_1/t_2)g(t_2)>0$, a contradiction. So $g$ has exactly one
zero in $[0,1]$, and it is $r(p)$. If $p_n\to p\in(0,1)$, every limit point of $r(p_n)$ is a zero of $F_p(t)-t$ by
continuity of $(p,t)\mapsto F_p(t)$, hence equals $r(p)$; so $r$ is continuous.
\end{proof}

\begin{remark}\label{rem:iteration}
Let $r_0=p$ and $r_{d+1}=F_p(r_d)$. By induction, $r_d$ is the probability that $x\leftrightarrow P(x)$ inside the part
of $H_x$ within $d$ generations below $x$, so $r_d$ increases to $r(p)$ by continuity of measure. This identifies $r(p)$
as the limit of the iteration started at $p$, without reference to the uniqueness argument.
\end{remark}

Fix vertices $u,v$ and a vertex $a$ with $u,v\in D(a)$, and let $S$ be the smallest subtree of $T$ containing $u$, $v$
and $a$. Since $a$ is an ancestor of $u$ and $v$, the set $S\setminus\{a\}$ is closed under $P$. Let
\[
f_{u,v,a}(p)=\PP_p\bigl(u\leftrightarrow v\text{ by an open path in the subgraph of }\GP\text{ induced by }D(a)\bigr).
\]

\begin{lemma}\label{lem:finite}
For all vertices $u,v$ and every vertex $a$ with $u,v\in D(a)$, the function $f_{u,v,a}$ is continuous on $(0,1)$.
\end{lemma}

\begin{proof}
The set $D(a)\setminus S$ is the disjoint union of the cones $D(z)$ over the vertices $z\notin S$ with $P(z)\in S$, and
there are finitely many of them. Fix such a $z$ and let $t=P(z)$. The edges between $D(z)$ and the rest of $D(a)$ are the
tree edge $\{z,t\}$, the edges $\{y,t\}$ for the children $y$ of $z$, and, when $t\neq a$, the edge $\{z,P(t)\}$; so
$D(z)$ is attached to $S$ only at $t$ and possibly at $s=P(t)\in S$. If $t=a$, the cone $D(z)$ is attached to $S$ at one
vertex and does not affect connections between vertices of $S$. If $t\neq a$, let $K_z$ be the graph $H_z$ together with
the vertex $s$ and the edge $\{z,s\}$. The only edge of $K_z$ at $s$ is $\{z,s\}$, so $K_z$ connects $s$ to $t$ if and
only if $\{z,s\}$ is open and $z\leftrightarrow t$ in $H_z$, an event of probability $pr(p)$.

The graphs $K_z$ use disjoint sets of edges, disjoint from the edges of $\GP$ inside $S$. Replace
each by one additional edge between $s$ and $t$, open exactly when the subgraph connects $s$ to $t$. A finite open path
from $u$ to $v$ in $D(a)$ becomes an open walk in the resulting finite multigraph on $S$ after each excursion into a cone
is replaced by its additional edge, and conversely each open additional edge can be expanded into a finite open path
through its cone. So $f_{u,v,a}(p)$ is the probability that $u$ and $v$ are connected in a finite multigraph whose edges
are open independently, with probability $p$ for the edges of $\GP$ inside $S$ and probability $pr(p)$ for the
additional edges. This is a polynomial in $p$ and $pr(p)$, hence continuous on $(0,1)$ by Lemma~\ref{lem:cone}.
\end{proof}

\section{Proof of Theorem~\ref{thm:main}}\label{sec:proof}

\begin{lemma}\label{lem:boundary}
For every vertex $x$, the set $B(x)$ of edges of $\GP$ with exactly one endpoint in $D(x)$ is
\[
B(x)=\bigl\{\{x,P(x)\},\{x,P^2(x)\}\bigr\}\cup\bigl\{\{y,P(x)\}: P(y)=x\bigr\},
\]
which has $b+2$ elements. For every vertex $w$, the sets $B(w),B(P^3(w)),B(P^6(w)),\dots$ are pairwise disjoint.
\end{lemma}

\begin{proof}
An edge of $\GP$ joins a vertex $y$ to $P(y)$ or to $P^2(y)$. It leaves $D(x)$ exactly when $y\in D(x)$ and the other
endpoint is not, which happens for $y=x$ (both edges) and for the children $y$ of $x$ (the edge to $P^2(y)=P(x)$). Every
edge of $B(x)$ has one endpoint at height $h(x)-1$ or $h(x)$. Every endpoint of an edge of $B(P^{3k}(w))$ has height at
least $h(w)+3k-1$, while every edge of $B(P^{3j}(w))$ has an endpoint at height at most $h(w)+3j$. For $j<k$ we have
$h(w)+3j<h(w)+3k-1$, so the two sets are disjoint.
\end{proof}

\begin{lemma}\label{lem:twopoint}
For all vertices $u,v$, the function $p\mapsto\PP_p(u\leftrightarrow v)$ is continuous on $(0,1)$.
\end{lemma}

\begin{proof}
Let $w$ be the nearest common ancestor of $u$ and $v$ in $T$, and let $a_n=P^n(w)$. If all edges of $B(a_k)$ are closed
for some $k\le n$, then the open cluster of $u$ is contained in $D(a_k)\subseteq D(a_n)$, and $u\leftrightarrow v$ in
$\GP$ if and only if $u\leftrightarrow v$ inside $D(a_n)$. The sets $B(a_k)$, $k\in\{0,3,6,\dots\}$, $k\le n$, are
$\lfloor n/3\rfloor+1$ fixed, pairwise disjoint sets of $b+2$ edges by Lemma~\ref{lem:boundary}, so
\begin{equation}\label{eq:tail}
0\le\PP_p(u\leftrightarrow v)-f_{u,v,a_n}(p)\le\bigl[1-(1-p)^{b+2}\bigr]^{\lfloor n/3\rfloor+1}.
\end{equation}
For $0<\alpha<\beta<1$ and $p\in[\alpha,\beta]$ the right side is at most $[1-(1-\beta)^{b+2}]^{\lfloor n/3\rfloor+1}$,
which tends to $0$ as $n\to\infty$. So $\PP_p(u\leftrightarrow v)$ is a uniform limit on $[\alpha,\beta]$ of the
functions $f_{u,v,a_n}$, which are continuous by Lemma~\ref{lem:finite}.
\end{proof}

The following is the observation of Häggström, Peres and Schonmann~\cite[p.~87]{HPS99}; it also appears as~\cite[Lemma~3.5]{BCRTT}. We include the
short proof.

\begin{lemma}\label{lem:sprinkle}
Let $G$ be a locally finite graph and $p\in(0,1)$. If $q\mapsto\PP_q(u\leftrightarrow v)$ is right continuous at $p$
for all vertices $u,v$, then $\PP_p$-almost surely any two distinct clusters are joined by finitely many edges.
\end{lemma}

\begin{proof}
Fix $u,v$ and let $A$ be the event that the clusters of $u$ and $v$ are distinct and joined by infinitely many edges.
Couple $\PP_p$ and $\PP_q$, $q>p$, in the standard monotone way. Given the configuration at level $p$, each closed edge
is open at level $q$ independently with probability $(q-p)/(1-p)>0$. On $A$ there are infinitely many closed edges
joining the two clusters, so almost surely one of them opens and $u\leftrightarrow v$ at level $q$. Hence
$\PP_q(u\leftrightarrow v)-\PP_p(u\leftrightarrow v)\ge\PP_p(A)$, and letting $q\downarrow p$ gives $\PP_p(A)=0$. Taking
the union over the countably many pairs $u,v$ proves the lemma.
\end{proof}

\begin{proof}[Proof of Theorem~\ref{thm:main}]
For $p\in(0,1)$ combine Lemmas~\ref{lem:twopoint} and~\ref{lem:sprinkle}. For $p=0$ every cluster is a single vertex
and touching sets have at most one vertex. For $p=1$ there is one cluster.
\end{proof}

\section{Remarks}\label{sec:remarks}

\begin{remark}
The argument uses two features of $\GP$: each descendant cone is attached to the rest of the graph at no more than two
vertices, and the edge boundaries of the cones along an ancestral line have bounded size. For general nonunimodular
transitive graphs cluster repulsion remains open~\cite[Question~1.1]{BCRTT}.
\end{remark}

\begin{remark}\label{rem:sims}
The repository contains exact checks and simulations. Exact enumeration in rational arithmetic confirms the recursion
of Remark~\ref{rem:iteration} on nine finite networks (the largest with $13$ edges and $8192$ configurations), the
replacement of attached cones by single edges on three cores, and Lemma~\ref{lem:boundary} for $b=2,3,4$, at
$p\in\{1/5,11/20,4/5\}$. Simulations on binary descendant cones of depths $10$, $12$ and $14$ at $p\in\{0.55,0.65,0.75\}$
with $12$ seeds each count the edges between distinct components with at least $64$ vertices; the largest number of
edges between one pair was $14$. A finite window can split one cluster of $\GP$ into several components, so these counts
illustrate the geometry and are not a test of Theorem~\ref{thm:main}.
\end{remark}

\section*{Acknowledgements and funding}
This work was done by multiple AI agents, Claude (Anthropic) and Codex (OpenAI), connected to Math
Brain~\cite{MathBrainRepo} and working together inside the Brain framework~\cite{BrainFramework}, a shared, persistent
memory in which the agents record claims, proofs, failed attempts and methods, retrieve one another's results,
coordinate who works on what, and audit each other's work. For more on the details behind the Brain and Math Brain, see
the Math Brain paper~\cite{MathBrainPaper}, the Math Brain repository~\cite{MathBrainRepo} and the Brain framework
repository~\cite{BrainFramework}.

The research presented in this paper was supported by the European Research Council (ERC) under the European Union's
Horizon Europe research and innovation programme (grant agreement No.~101041711), by the Simons Foundation as part of
the Collaboration on the Mathematical and Scientific Foundations of Deep Learning, by Heights Labs, by Convex Nexus
Capital, by the Israel Science Foundation (grant number 2258/19), and by the Israel Science Foundation (ISF Grant
4101/25). Computational resources were provided in part by Google Cloud through its education credits programme (credit
530550103).

\section*{Declaration of competing interest}
The author holds equity in Heights Labs and in Convex Nexus Capital and has a management role at Convex Nexus Capital.
These entities are included in the funding declaration above.

\section*{Declaration of generative AI and AI-assisted technologies}
A Codex agent (OpenAI) proved Theorem~\ref{thm:main} and wrote the code of Remark~\ref{rem:sims}; a second Codex agent
checked the proof and the literature as a referee; Claude agents (Anthropic) checked the steps and prepared the
manuscript. The author directed the work and is responsible for the content of the paper and for decisions concerning
its dissemination.

\section*{Code availability}
The code and data of Remark~\ref{rem:sims} are in the repository
\begin{center}\url{https://github.com/yspennstate/grandparent-graph-cluster-repulsion}.\end{center}
They use Python~3 and its standard library only.

\appendix
\section{Notation}\label{app:notation}

\begin{center}
\begin{tabular}{p{3.0cm}p{8.0cm}p{3.1cm}}
\toprule
Symbol & Meaning & First appears\\
\midrule
$G=(V,E)$ & a transitive graph & Section~\ref{sec:intro}\\
$\PP_p$ & Bernoulli bond percolation with parameter $p$ & Section~\ref{sec:intro}\\
$\tau(C,D)$ & touching set of two clusters & Section~\ref{sec:intro}\\
$u\leftrightarrow v$ & $u$ and $v$ are joined by an open path & Section~\ref{sec:intro}\\
$\Gamma$, $\Gamma_x$, $m$ & automorphism group, stabilizer, Haar measure & Section~\ref{sec:intro}\\
$b$, $T$, $\xi$ & branching number, $(b+1)$-regular tree, fixed end & Section~\ref{sec:intro}\\
$P(x)$ & neighbour of $x$ towards $\xi$ & Section~\ref{sec:intro}\\
$\GP$ & grandparent graph & Section~\ref{sec:intro}\\
$h$ & height function, $h(P(x))=h(x)+1$ & Section~\ref{sec:structure}\\
$D(x)$ & descendant cone of $x$ & Section~\ref{sec:structure}\\
$H_x$ & subgraph induced by $D(x)\cup\{P(x)\}$ & Section~\ref{sec:cone}\\
$r(p)$, $F_p$ & cone reliability and its equation & \eqref{eq:cone}\\
$g$ & $F_p(t)-t$ & Lemma~\ref{lem:cone}\\
$r_d$ & reliability of a cone cut at depth $d$ & Remark~\ref{rem:iteration}\\
$S$ & tree spanned by $u$, $v$, $a$ & Section~\ref{sec:cone}\\
$f_{u,v,a}$ & connection probability inside $D(a)$ & Section~\ref{sec:cone}\\
$z$, $t$, $s$ & top of an attached cone and its attachment vertices & Lemma~\ref{lem:finite}\\
$K_z$ & $H_z$ with the vertex $s$ and the edge $\{z,s\}$ & Lemma~\ref{lem:finite}\\
$B(x)$ & edge boundary of $D(x)$ & Lemma~\ref{lem:boundary}\\
$w$, $a_n$ & nearest common ancestor, $P^n(w)$ & Lemma~\ref{lem:twopoint}\\
$\alpha$, $\beta$ & ends of a compact interval of parameters & Lemma~\ref{lem:twopoint}\\
$q$ & a second parameter, $q>p$ & Lemma~\ref{lem:sprinkle}\\
$A$ & event of infinitely many edges between two clusters & Lemma~\ref{lem:sprinkle}\\
\bottomrule
\end{tabular}
\end{center}

\begin{thebibliography}{99}
\bibitem{BCRTT} S.~Bell, T.~Chu, O.~Rodgers, G.~Terlov and A.~Tserunyan, Heavy repulsion of clusters in Bernoulli
percolation, preprint (2025), \arxiv{2509.10631}.
\bibitem{HJ06} O.~Häggström and J.~Jonasson, Uniqueness and non-uniqueness in percolation theory, \emph{Probab.\ Surv.}
\textbf{3} (2006), 289--344. \doi{10.1214/154957806000000096}
\bibitem{HPS99} O.~Häggström, Y.~Peres and R.~H.~Schonmann, Percolation on transitive graphs as a coalescent process:
relentless merging followed by simultaneous uniqueness, in \emph{Perplexing Problems in Probability}, Progr.\ Probab.\
\textbf{44}, Birkhäuser, Boston, 1999, 69--90.
\bibitem{Hut20} T.~Hutchcroft, Nonuniqueness and mean-field criticality for percolation on nonunimodular transitive
graphs, \emph{J.\ Amer.\ Math.\ Soc.} \textbf{33} (2020), no.~4, 1101--1165. \doi{10.1090/jams/953}
\bibitem{MathBrainPaper} Y.~Shmalo, \emph{Math Brain: shared, auditable memory for multi-agent mathematical research},
October 2026, \url{https://github.com/yspennstate/math-brain-public/tree/main/papers/math-brain-2026-10-24}.
\bibitem{MathBrainRepo} Math Brain repository, 2026, \url{https://github.com/yspennstate/math-brain-public}.
\bibitem{BrainFramework} Y.~Shmalo, Brain Framework, public software repository, 2026,
\url{https://github.com/yspennstate/brain-framework-public}.
\bibitem{MS26} L.~Ménard and D.~Sénizergues, Percolation on grandparent graphs of Bienaymé trees, preprint (2026),
\arxiv{2611.09762}.
\bibitem{Spa09} I.~Špakulová, Critical percolation of virtually free groups and other tree-like graphs, \emph{Ann.\
Probab.} \textbf{37} (2009), no.~6, 2262--2296. \doi{10.1214/09-AOP458}
\bibitem{Tang19} P.~Tang, Heavy Bernoulli-percolation clusters are indistinguishable, \emph{Ann.\ Probab.} \textbf{47}
(2019), no.~6, 4077--4115. \doi{10.1214/19-AOP1354}
\bibitem{Tim06a} Á.~Timár, Neighboring clusters in Bernoulli percolation, \emph{Ann.\ Probab.} \textbf{34} (2006),
no.~6, 2332--2343. \doi{10.1214/009117906000000485}
\bibitem{Tim06b} Á.~Timár, Percolation on nonunimodular transitive graphs, \emph{Ann.\ Probab.} \textbf{34} (2006),
no.~6, 2344--2364. \doi{10.1214/009117906000000494}
\end{thebibliography}
\end{document}
